\chapter{Азбука Морзе: что мы знаем о языке, на котором говорят \nt{GLUT}- и \nt{GABA}-нейроны}
\chaptermark{Азбука Морзе}
\label{sec:code}

\section{Алфавит из одного знака}

В азбуке Морзе два знака, точка и тире, и паузы трёх длин между ними. У нейрона, как мы видели в главе~\ref{sec:neurontypes}, алфавит ещё беднее: знак один. Импульс длится около миллисекунды, и все импульсы одного нейрона одинаковы по высоте и форме. Поэтому всё сообщение --- в паузах, то есть в последовательности моментов $t_1,t_2,t_3,\dots$ Это азбука без тире, в которой смысл несёт только ритм точек (рис.~\ref{fig:morse}).

\begin{figure}[t]
\centering
\begin{tikzpicture}[>=Stealth,x=0.08cm,y=1cm]
% Morse letter: dot, dash, dot
\node[left,font=\small] at (-6,2.55) {Морзе:};
\fill[black] (0,2.45) rectangle (6,2.65);
\fill[black] (14,2.45) rectangle (32,2.65);
\fill[black] (40,2.45) rectangle (46,2.65);
\node[right,font=\small,gray!40!black] at (52,2.55) {точка, тире, точка --- смысл в длительностях и паузах};
% stimulus onset
\node[left,font=\small] at (-6,0.4) {нейрон:};
\draw[->,thick,green!45!black] (0,-0.55) -- (0,-0.05);
\node[below,font=\scriptsize,green!45!black] at (0,-0.55) {раздражитель};
\draw[gray!70!black] (0,0) -- (152,0) node[right,black] {\small $t$};
% spikes: single ones blue, the burst red
\foreach \s in {14,26,41,88,112,131}{\draw[very thick,blue!60!black] (\s,0) -- (\s,0.8);}
\foreach \s in {60,63,66,69}{\draw[very thick,red!70!black] (\s,0) -- (\s,0.8);}
% latency
\draw[<->,thick] (0,1.1) -- node[above,font=\scriptsize] {$t_1$: время первого импульса} (14,1.1);
% burst
\draw[decorate,decoration={brace,amplitude=4pt},red!70!black] (58,0.95) -- (71,0.95) node[midway,above=4pt,font=\scriptsize] {пачка --- «тире»};
% counting windows
\foreach \a/\b/\n in {0/50/3,50/100/5,100/150/2}{
  \draw[decorate,decoration={brace,amplitude=4pt,mirror},gray!50!black] (\a+1,-0.2) -- (\b-1,-0.2) node[midway,below=4pt,font=\small] {$n=\n$};}
\node[font=\scriptsize,gray!40!black] at (75,-1.05) {частотный код: число импульсов $n$ в каждом окне};
\foreach \x in {0,50,100,150}{\draw[gray!60!black] (\x,-0.06) -- (\x,0.06);}
\end{tikzpicture}
\caption{Азбука Морзе и азбука нейрона. Одну и ту же последовательность импульсов можно прочитать по-разному: сосчитать импульсы в окнах по $50\ms$ (раздел~\ref{sec:rate}), измерить задержку $t_1$~\eqref{eq:latency} или заметить пачку --- «тире»~\eqref{eq:burst}.}
\label{fig:morse}
\end{figure}

Какие паузы вообще возможны? Снизу их ограничивает невосприимчивость: после импульса нейрон около миллисекунды не может сработать снова, так что больше нескольких сотен импульсов в секунду не бывает. Сверху не ограничивает ничто: нейрон может молчать минутами. У \nt{GLUT}-нейронов коры частоты лежат в пределах от долей импульса до нескольких десятков импульсов в секунду, и большинство нейронов большую часть времени работает у нижней границы.

В главах~\ref{sec:glutcomputer} и~\ref{sec:gaba} мы попутно принимали то один, то другой способ записи чисел импульсами. Теперь спросим прямо: на каком языке \nt{GLUT}- и \nt{GABA}-нейроны говорят друг с другом? Полного ответа нет, этот язык не расшифрован. Но кое-что о нём известно твёрдо, и почти всё известное можно получить, рассуждая как инженер связи: ёмкость канала, шум, приёмник, цена передачи.

\section{Сколько можно передать}

Начнём с верхней границы. Пусть получатель различает моменты импульсов с точностью $\Delta t$: сдвиги меньше $\Delta t$ для него неразличимы. Разрежем время на клетки длины $\Delta t$, меньшие времени невосприимчивости, так что в каждой клетке либо один импульс, либо ни одного (рис.~\ref{fig:cells}). При средней частоте $r$ импульс попадает в клетку с вероятностью $p=r\Delta t$. Больше всего информации поток несёт, когда клетки независимы, и тогда каждая клетка даёт энтропию $-p\log_2p-(1-p)\log_2(1-p)\approx p\log_2(e/p)$ бит при $p\ll1$. Разделив на $\Delta t$, получим предельную скорость передачи и её долю на один импульс~\cite{MacKay1952}:
\begin{empheq}[box=\keybox]{equation}
R_{\max} \;\approx\; r\,\log_2\frac{e}{r\,\Delta t}\ \ \text{бит/с},
\qquad
\frac{R_{\max}}{r} \;\approx\; \log_2\frac{e}{r\,\Delta t}\ \ \text{бит на импульс}.
\label{eq:spikeentropy}
\end{empheq}
При $r=10$ импульсов в секунду и $\Delta t=1\ms$ это около восьми бит на импульс и восьмидесяти бит в секунду.

\begin{figure}[t]
\centering
\begin{tikzpicture}[>=Stealth,x=0.5cm,y=1cm]
% time cut into cells of width Delta t; a cell holds one impulse or none
\foreach \k in {0,...,23}{\draw[gray!55] (\k,0) rectangle (\k+1,0.7);}
\foreach \k in {2,7,8,15,21}{
  \fill[red!10] (\k,0) rectangle (\k+1,0.7);
  \draw[very thick,red!70!black] (\k+0.5,0.05) -- (\k+0.5,0.65);}
\foreach \k in {0,...,23}{
  \pgfmathtruncatemacro{\bit}{(\k==2)||(\k==7)||(\k==8)||(\k==15)||(\k==21)}
  \ifnum\bit=1 \def\bitcol{red!70!black}\else \def\bitcol{gray!60!black}\fi
  \node[font=\scriptsize,text=\bitcol] at (\k+0.5,-0.25) {\bit};}
\draw[<->,gray!60!black] (0,0.95) -- (1,0.95) node[midway,above,font=\scriptsize,black] {$\Delta t$};
\node[font=\small,anchor=east] at (-0.3,0.35) {импульсы};
\node[font=\small,anchor=east] at (-0.3,-0.25) {биты};
\draw[decorate,decoration={brace,amplitude=5pt,mirror},gray!60!black] (0,-0.55) -- (24,-0.55)
  node[midway,below=6pt,font=\scriptsize,black,align=center] {получатель, который только считает: «$n=5$»;\\ получатель, который видит клетки: какие именно $5$ из $24$ --- это $\log_2\binom{24}{5}\approx15$ бит};
\end{tikzpicture}
\caption{Откуда берётся ёмкость~\eqref{eq:spikeentropy}. Время нарезано на клетки длины $\Delta t$, в каждой клетке импульс либо есть ($1$), либо нет ($0$): поток импульсов --- это двоичная строка. Импульс попадает в клетку с вероятностью $p=r\Delta t$. Чем мельче клетки, тем длиннее строка и тем больше бит; счётчик импульсов теряет почти всё, что записано в том, \emph{где} стоят единицы.}
\label{fig:cells}
\end{figure}

Биты на импульс зависят от точности времени лишь логарифмически: при точности $10\ms$ останется около пяти бит на импульс. Но если получатель только считает импульсы за секунду, то при $r=10$ он получит меньше двух бит за всю секунду (раздел~\ref{sec:rate}).

\begin{keyblock}
\begin{proposition}[Ёмкость импульсного канала]
\label{prop:capacity}
Поток импульсов со средней частотой $r$, прочитанный с точностью времени $\Delta t$, несёт не больше~\eqref{eq:spikeentropy}. Почти вся эта ёмкость сидит во \emph{времени} импульсов: получатель, который только считает импульсы, извлекает из того же потока в десятки раз меньше, чем тот, кто различает их моменты с точностью до миллисекунды.
\end{proposition}
\end{keyblock}

Это верхняя граница, а не измерение. Измерения на чувствительных нейронах, для которых известен раздражитель, дают от одного до нескольких бит на импульс; в лучших случаях это до половины границы, вычисленной для их точности~\cite{Rieke1997}. Значит, по крайней мере на входе в нервную систему время импульсов используется, а не пропадает.

\section{Шумный канал}

Ёмкость~\eqref{eq:spikeentropy} посчитана без шума. О том, где он возникает, твёрдо известны три факта.

\emph{Сам генератор импульсов точен.} Если в нейрон коры, вынутый из мозга вместе с кусочком ткани, много раз подать один и тот же ток, быстро меняющийся во времени, импульсы повторяются с точностью около миллисекунды~\cite{Mainen1995}. Если ток постоянный, моменты импульсов от раза к разу расползаются: точное время порождают быстрые перепады входа, а не сам нейрон.

\emph{Синапс ненадёжен.} Импульс, добежавший до \nt{GLUT}-синапса, выбрасывает порцию медиатора лишь с некоторой вероятностью $p$. На синапсах гиппокампа больше половины импульсов до получателя просто не доходят, и причина --- именно случайность выброса, а не сбои проведения~\cite{AllenStevens1994}. Для связиста это канал со стираниями; несколько синапсов между двумя нейронами отчасти спасают дело.

\emph{В живой коре поток импульсов выглядит случайным.} Интервалы между импульсами разбросаны примерно так же, как у случайного пуассоновского потока, а при повторении одного и того же раздражителя число импульсов за окно меняется так, что его дисперсия близка к среднему~\cite{Softky1993}.

Как точный элемент порождает случайный поток? Нейрон получает тысячи входов, каждый ненадёжен, а возбуждение и торможение уравновешены, как в главе~\ref{sec:gaba}. Напряжение мембраны дрожит у порога, и моменты импульсов задают эти дрожания. Шум ли это или сообщение, которое мы не умеем прочитать, --- во многом открытый вопрос. Часть «шума» уже оказалась сообщением: в зрительной коре значительная доля разброса следует за движениями самого животного~\cite{Stringer2019}. Трудность здесь принципиальная: хорошо сжатое сообщение для того, кто не знает кода, неотличимо от случайного потока.

\section{Частота: медленно, но надёжно}
\label{sec:rate}

Самый простой код --- \emph{частотный}: число записано в том, сколько импульсов нейрон выдаёт за окно $T$. Ни стирания, ни дрожание моментов такому коду не страшны. Но у него есть цена. Если поток пуассоновский, число импульсов $n$ за окно имеет среднее $rT$ и разброс $\sqrt{rT}$:
\begin{empheq}[box=\keybox]{equation}
\frac{\sigma_n}{\bar n} \;=\; \frac{1}{\sqrt{r\,T}} .
\label{eq:rateerror}
\end{empheq}
Чтобы узнать частоту с точностью в десять процентов, нужно сто импульсов, при двадцати импульсах в секунду --- пять секунд. Соседние уровни различимы, если они отстоят на пару разбросов, поэтому до частоты $r$ за время $T$ различается около $\sqrt{rT}$ уровней: при $r=10$ и $T=1\,\text{с}$ это три-четыре уровня, меньше двух бит.
Мозг так медленно не работает. Человек узнаёт животное на картинке, показанной на мгновение, и ответ в мозге появляется примерно через $150\ms$~\cite{Thorpe1996}. По грубой оценке, за это время сигнал проходит около десятка последовательных ступеней обработки, по $10$--$20\ms$ на ступень. При обычных частотах коры каждый нейрон успевает выдать на ступени ноль или один импульс, и его частота в таком окне не определена. Выходов два: взять много нейронов или смотреть на время.

\emph{Много нейронов.} Пусть $N$ нейронов несут одно и то же число, и получатель складывает их импульсы. Если их шумы независимы, в~\eqref{eq:rateerror} вместо $rT$ встаёт $NrT$: сто нейронов при $r=20$ за $15\ms$ дают тридцать импульсов и точность около двадцати процентов. Пространство заменяет время. Однако шумы соседних нейронов коры не независимы: коэффициент корреляции чисел импульсов у пары соседей обычно около $0.1$~\cite{Zohary1994}. Если он равен $c$, дисперсия среднего по $N$ нейронам равна
\begin{empheq}[box=\keybox]{equation}
\sigma^2_N \;=\; \sigma^2_1\,\frac{1+(N-1)\,c}{N}
\;\xrightarrow[N\to\infty]{}\; c\,\sigma^2_1 .
\label{eq:correlated}
\end{empheq}

\begin{keyblock}
\begin{proposition}[Предел усреднения]
\label{prop:pooling}
При коррелированном шуме усреднение по группе уменьшает ошибку не больше чем в $1/\sqrt{c}$ раз. При $c=0.1$ это втрое, и такой выигрыш достигается уже на нескольких десятках нейронов; добавлять больше бесполезно.
\end{proposition}
\end{keyblock}

Не всякая корреляция одинаково вредна. Ограничивает точность только та часть общего шума, которая \emph{похожа на сам сигнал}, то есть сдвигает всю группу так же, как её сдвинуло бы изменение измеряемой величины~\cite{MorenoBote2014}. Сколько такого шума в мозге на самом деле, ещё выясняют.

Есть и другой способ использовать много нейронов: записать число в том, \emph{какой} нейрон сработал, как в определителе направления на звук (рис.~\ref{fig:itd}). Это код \emph{местом}, самый надёжный из известных: у чувствительных нейронов номер провода говорит, какая точка кожи тронута или какой высоты звук, и это установлено многократно. Он дорог по числу нейронов, зато быстр: на сообщение уходит одна задержка.

\section{Время: быстро, но от чего отсчитывать}

Второй выход --- писать число во время импульса. Возьмём нейрон~\eqref{eq:glutneuron} и подадим на него с момента $t=0$ постоянный вход, который поднял бы напряжение до уровня $U$. Напряжение растёт как $V=U\bigl(1-e^{-t/\tau}\bigr)$ и достигает порога $\theta$ в момент
\begin{empheq}[box=\keybox]{equation}
t_1 \;=\; \tau\,\ln\frac{U}{U-\theta}, \qquad U>\theta .
\label{eq:latency}
\end{empheq}
Сильный вход --- ранний импульс, слабый --- поздний, подпороговый --- никакого. При $\tau=10\ms$ вход $U=4\theta$ даёт первый импульс через $2.9\ms$, $U=2\theta$ --- через $6.9\ms$, $U=1.2\theta$ --- через $17.9\ms$. Один-единственный импульс несёт число.

Но от какого момента отсчитывать $t_1$? Часов нет, и получатель не знает, когда начался раздражитель. Есть три ответа.

\emph{Относительная задержка.} Два нейрона видят один и тот же раздражитель, и разность их задержек $t_1^A-t_1^B$ зависит только от их входов, а не от момента начала. Сравнивают моменты детекторы совпадений с задержками (рис.~\ref{fig:itd}). Если $N$ нейронов выдали по одному импульсу, сам \emph{порядок} их прихода несёт до $\log_2N!$ бит: для десяти нейронов около двадцати двух бит, тогда как ответ «сработал или нет» дал бы не больше десяти. В сетчатке показано, что по относительным задержкам первых импульсов её выходных нейронов картинку можно восстановить быстро и надёжно~\cite{Gollisch2008}.

\emph{Залп как начало слова.} Резкое изменение раздражителя заставляет сработать почти одновременно множество нейронов. Такой залп сам служит отметкой начала, как длинная пауза между словами в азбуке Морзе, и получатель отсчитывает задержки от него.

\emph{Местный ритм.} В главе~\ref{sec:gaba} петля \nt{GLUT}--\nt{GABA} порождала местный ритм с периодом $12$--$50\ms$. Это не часы, но внутри его цикла нейрон с сильным входом успевает сработать раньше, чем нейрон со слабым, и~\eqref{eq:latency} превращается в код \emph{фазой}: число записано в том, в какой части цикла пришёл импульс. Такой код наблюдали: нейроны, отвечающие за определённое место в пространстве, срабатывают всё раньше в цикле медленного местного ритма по мере того, как животное проходит через «своё» место~\cite{OKeefe1993}. Насколько широко мозг пользуется фазой, пока неизвестно.

\section{Код выбирает получатель}

В последовательности импульсов есть и частота, и времена, и порядок. Что из этого будет прочитано, решает получатель, и решает одним параметром --- своей постоянной времени $\tau$.

Усредним уравнение~\eqref{eq:glutneuron} по времени. Если на нейрон приходят независимые потоки с частотами $r_j$, то среднее напряжение и его относительное дрожание равны
\begin{empheq}[box=\keybox]{equation}
\bar V \;=\; \tau\sum_j w_j\,r_j ,
\qquad
\frac{\sigma_V}{\bar V} \;\approx\; \frac{1}{\sqrt{2\,r_{\text{вх}}\,\tau}} ,
\label{eq:reader}
\end{empheq}
где второе равенство написано для одинаковых весов, а $r_{\text{вх}}$ --- суммарная частота всех входов. Когда $\tau$ велика, напряжение почти не дрожит и повторяет взвешенную сумму частот: получатель --- \emph{интегратор}, он читает частоты и забывает времена. Когда $\tau$ мала, напряжение успевает упасть между импульсами, и порог достигается только тогда, когда несколько импульсов пришли в пределах окна~\eqref{eq:window}: получатель --- \emph{детектор совпадений}, он читает времена (рис.~\ref{fig:reader}). Тысяча входов по пять импульсов в секунду при $\tau=10\ms$ дают дрожание около десяти процентов, почти чистое интегрирование. Если торможение, как в главе~\ref{sec:gaba}, укоротило $\tau$ до $2\ms$, дрожание вырастает до двадцати с лишним процентов, и нейрон начинает замечать совпадения.

\begin{figure}[t]
\centering
\begin{tikzpicture}[>=Stealth,x=0.115cm,y=1cm]
% common input: the same spike train for both receivers
\foreach \s in {6,21,22,38,52,55.5,59,62.5,66,69.5,73,76.5,80,83.5,87,90.5,94,97.5}{\draw[thick,gray!40!black] (\s,5.25) -- (\s,5.65);}
\draw[gray!70!black] (0,5.25) -- (105,5.25);
\node[left,font=\small] at (-1,5.45) {вход};
\node[above,font=\scriptsize,gray!40!black] at (13.5,5.65) {редко};
\node[above,font=\scriptsize,gray!40!black] at (21.5,6.0) {пара};
\node[above,font=\scriptsize,gray!40!black] at (75,5.65) {часто};
% axes and thresholds
\foreach \y/\th in {2.7/2.5,0/1.5}{
  \draw[gray!70!black] (0,\y) -- (106,\y) node[right,black] {\small $t$};
  \draw[densely dashed,gray!60!black] (0,\y+0.6*\th) -- (104,\y+0.6*\th) node[right,black,font=\small] {$\theta$};
}
\node[anchor=south west,font=\small] at (0,4.3) {$\tau=10\ms$: интегратор};
\node[anchor=south east,font=\small] at (104,1.0) {$\tau=2\ms$: детектор совпадений};
% integrator: tau=10 ms, theta=2.5w
\begin{scope}[yshift=2.7cm]
\foreach \a/\b/\v in {0/6/0.000,6/21/1.000,21/22/1.223,22/38/2.107,38/52/1.425,52/55.5/1.351,55.5/59/1.952,59/62.5/2.376,62.5/66/0.000,66/69.5/1.000,69.5/73/1.705,73/76.5/2.201,76.5/80/0.000,80/83.5/1.000,83.5/87/1.705,87/90.5/2.201,90.5/94/0.000,94/97.5/1.000,97.5/104/1.705}{\draw[very thick,blue!60!black] plot[domain=\a:\b,samples=14] (\x,{0.6*\v*exp(-(\x-\a)/10)});}
\foreach \s/\p/\q in {6/0.000/1.000,21/0.223/1.223,22/1.107/2.107,38/0.425/1.425,52/0.351/1.351,55.5/0.952/1.952,59/1.376/2.376,62.5/1.674/2.500,62.5/2.500/0.000,66/0.000/1.000,69.5/0.705/1.705,73/1.201/2.201,76.5/1.551/2.500,76.5/2.500/0.000,80/0.000/1.000,83.5/0.705/1.705,87/1.201/2.201,90.5/1.551/2.500,90.5/2.500/0.000,94/0.000/1.000,97.5/0.705/1.705}{\draw[very thick,blue!60!black] (\s,0.6*\p) -- (\s,0.6*\q);}
\foreach \s in {62.5,76.5,90.5}{\draw[->,thick,red!70!black] (\s,1.58) -- (\s,1.95);}
\end{scope}
% coincidence: tau=2 ms, theta=1.5w
\begin{scope}[yshift=0cm]
\foreach \a/\b/\v in {0/6/0.000,6/21/1.000,21/22/1.001,22/38/0.000,38/52/1.000,52/55.5/1.001,55.5/59/1.174,59/62.5/1.204,62.5/66/1.209,66/69.5/1.210,69.5/73/1.210,73/76.5/1.210,76.5/80/1.210,80/83.5/1.210,83.5/87/1.210,87/90.5/1.210,90.5/94/1.210,94/97.5/1.210,97.5/104/1.210}{\draw[very thick,blue!60!black] plot[domain=\a:\b,samples=14] (\x,{0.6*\v*exp(-(\x-\a)/2)});}
\foreach \s/\p/\q in {6/0.000/1.000,21/0.001/1.001,22/0.607/1.500,22/1.500/0.000,38/0.000/1.000,52/0.001/1.001,55.5/0.174/1.174,59/0.204/1.204,62.5/0.209/1.209,66/0.210/1.210,69.5/0.210/1.210,73/0.210/1.210,76.5/0.210/1.210,80/0.210/1.210,83.5/0.210/1.210,87/0.210/1.210,90.5/0.210/1.210,94/0.210/1.210,97.5/0.210/1.210}{\draw[very thick,blue!60!black] (\s,0.6*\p) -- (\s,0.6*\q);}
\foreach \s in {22}{\draw[->,thick,red!70!black] (\s,0.98) -- (\s,1.35);}
\end{scope}
\foreach \x in {0,50,100}{\draw[gray!60!black] (\x,-0.06) -- (\x,0.06) node[below,font=\scriptsize] {\x\,мс};}
\end{tikzpicture}
\caption{Один и тот же вход --- два разных получателя. Каждый импульс поднимает напряжение на $w$, после чего оно стекает, как в модели нейрона~\eqref{eq:glutneuron}; красные стрелки --- импульсы получателя. Медленный получатель ($\tau=10\ms$, $\theta=2.5w$) срабатывает там, где импульсов \emph{много}, и не замечает пары. Быстрый ($\tau=2\ms$, $\theta=1.5w$) срабатывает только на пару в пределах окна~\eqref{eq:window}, а частый, но не совпадающий поток пропускает.}
\label{fig:reader}
\end{figure}

Измерения показывают, что в активной коре проводимость мембраны возрастает в несколько раз по сравнению с покоем~\cite{Destexhe2003}. Значит, $\tau$ там короче, и нейроны коры в рабочем режиме ближе к детекторам совпадений, чем к интеграторам.

\begin{keyblock}
\begin{proposition}[Код задаёт получатель]
\label{prop:reader}
Одна и та же последовательность импульсов для медленного получателя несёт частоту, для быстрого --- времена, для объёма, в который выброшен медиатор, --- уровень концентрации, сглаженный за секунды (глава~\ref{sec:serocomputer}). Какой код используется, определяется не отправителем, а постоянной времени получателя, и её подстраивает торможение.
\end{proposition}
\end{keyblock}
\section{Точки и тире: синапс как фильтр}

До сих пор синапс был у нас проводом с весом. На самом деле его ответ зависит от предыдущих импульсов, и это даёт языку нейронов второй знак.

\emph{Истощение: синапс передаёт изменения.} Каждый выброс тратит долю $U$ запаса медиатора, готового к выбросу~\cite{Tsodyks1997}, а запас восстанавливается с постоянной времени $\tau_{\text{вос}}$ порядка сотен миллисекунд. При частоте $r$ амплитуда ответа на импульс устанавливается приблизительно на уровне $A(r)=A_0/(1+U r\tau_{\text{вос}})$, где $A_0$ --- ответ отдохнувшего синапса. Ток, который получатель получает в единицу времени, равен
\begin{empheq}[box=\keybox]{equation}
r\,A(r) \;=\; \frac{A_0\,r}{1+U r\,\tau_{\text{вос}}}
\;\xrightarrow[r\to\infty]{}\; \frac{A_0}{U\tau_{\text{вос}}} .
\label{eq:depression}
\end{empheq}
При $U=0.5$ и $\tau_{\text{вос}}=0.8\,\text{с}$ насыщение начинается уже с $1/(U\tau_{\text{вос}})=2.5$ импульса в секунду. Если частота выросла с пяти до двадцати, установившийся ток вырастет всего на треть. Зато первые импульсы на новой частоте приходят, пока запас ещё прежний, и ток на мгновение подскакивает вчетверо, а затем за $\tau_{\text{вос}}/(1+Ur\tau_{\text{вос}})\approx90\ms$ спадает (рис.~\ref{fig:depression}).

Такой синапс передаёт не частоту, а её \emph{изменение}, по существу --- относительное изменение $\Delta r/r$~\cite{Abbott1997depression}. Для электронщика это фильтр верхних частот, поставленный прямо в провод.

\begin{figure}[t]
\centering
\begin{tikzpicture}[>=Stealth,x=12.5cm,y=1cm]
% input rate
\draw[->,gray!70!black] (0,2.6) -- (0.9,2.6) node[right,black] {\small $t$};
\draw[->,gray!70!black] (0,2.6) -- (0,4.3) node[above,black] {\small $r$};
\draw[very thick,blue!60!black] (0,2.9) -- (0.2,2.9) -- (0.2,3.8) -- (0.8,3.8);
\node[left,font=\scriptsize] at (0,2.9) {$5$};
\node[left,font=\scriptsize] at (0,3.8) {$20$};
\node[right,font=\small,blue!60!black] at (0.4,4.1) {частота на входе синапса};
% output current
\draw[->,gray!70!black] (0,0) -- (0.9,0) node[right,black] {\small $t$};
\draw[->,gray!70!black] (0,0) -- (0,2.45) node[left,black] {\small $rA$};
\draw[densely dashed,gray!60!black] (0,0.667) -- (0.8,0.667);
\draw[very thick,red!70!black] (0,0.5) -- (0.2,0.5) -- (0.2,2.0)
   plot[domain=0.2:0.8,samples=60] (\x,{0.667+(2.0-0.667)*exp(-(\x-0.2)/0.089)});
\node[left,font=\scriptsize] at (0,0.5) {$1.7$};
\node[left,font=\scriptsize] at (0,2.0) {$6.7$};
\node[right,font=\scriptsize] at (0.8,0.667) {$2.2$};
\node[right,font=\small,red!70!black] at (0.28,1.6) {ток получателю: подскок и спад за $\approx90\ms$};
\draw[gray!60!black] (0.2,-0.08) -- (0.2,0.08); \node[below,font=\scriptsize] at (0.2,0) {$0.2\,\text{с}$};
\draw[gray!60!black] (0.8,-0.08) -- (0.8,0.08); \node[below,font=\scriptsize] at (0.8,0) {$0.8\,\text{с}$};
\end{tikzpicture}
\caption{Истощающийся синапс по формуле~\eqref{eq:depression} при $U=0.5$, $\tau_{\text{вос}}=0.8\,\text{с}$; ток в единицах $A_0$ за секунду, штриховая линия --- установившийся ток: он вырос лишь с $1.7$ до $2.2$, а в момент скачка --- до $6.7$. Синапс сообщает получателю о том, что частота изменилась, а не о том, какая она.}
\label{fig:depression}
\end{figure}

\emph{Облегчение: пачка как тире.} У других синапсов вероятность выброса мала, но после каждого импульса на десятки миллисекунд растёт. Одиночный импульс через такой синапс чаще всего теряется. А \emph{пачка} --- несколько импульсов подряд с интервалами в несколько миллисекунд --- проходит почти наверняка. Даже без облегчения вероятность, что пропадут все $k$ импульсов пачки, равна
\begin{empheq}[box=\keybox]{equation}
P_{\text{потери}} \;=\; (1-p)^k ,
\label{eq:burst}
\end{empheq}
при $p=0.2$ это $0.8$ для одиночного импульса и $0.41$ для пачки из четырёх; облегчение уменьшает её ещё сильнее. Так у нейрона появляется второй знак: одиночный импульс --- точка, пачка --- тире. Истощающийся синапс лучше всего передаёт первый импульс после паузы; облегчающийся пропускает пачки и глушит одиночные точки. То, что пачки передаются надёжнее, измерено; что мозг действительно пользуется ими как отдельным знаком, --- правдоподобная гипотеза~\cite{Lisman1997}.

\section{Как говорят \nt{GABA}-нейроны}

Самые многочисленные из них в коре --- быстрые~\cite{Tremblay2016}. Их импульсы короче, чем у \nt{GLUT}-нейронов, они могут выдавать их ровно и подолгу с частотой в сотни в секунду, почти не уставая. Их синапсы надёжнее: связь с каждым получателем состоит из многих мест выброса, и импульс почти никогда не теряется. Их выходы садятся рядом с тем местом получателя, где рождается его импульс, так что они управляют не тем, как получатель складывает входы, а тем, сработает ли он и когда.

Сами они получают возбуждение от многих соседних \nt{GLUT}-нейронов и сообщают о суммарной активности вокруг --- ровно то, что нужно для деления из главы~\ref{sec:gaba}. Наконец, быстрые \nt{GABA}-нейроны связаны друг с другом и легко срабатывают вместе; именно они задают местный ритм, о котором шла речь выше~\cite{Cardin2009}.

На языке азбуки Морзе \nt{GABA}-нейроны отвечают не за буквы, а за \emph{паузы}. Они режут непрерывный поток \nt{GLUT}-импульсов на окна, внутри которых получатель может сработать, сужают окна совпадения и приглушают весь поток, когда он становится слишком громким.

\begin{keyblock}
\begin{proposition}[Разделение ролей]
\label{prop:punctuation}
\nt{GLUT}-нейроны несут содержание сообщения: \emph{что} и \emph{сколько}. Быстрые \nt{GABA}-нейроны несут его разметку: \emph{когда} можно говорить и \emph{насколько громко}; ещё один класс, тормозя тормозящих, решает, \emph{для кого} ворота открыты. Отдельные части этого разделения измерены; как общее правило это рабочая гипотеза.
\end{proposition}
\end{keyblock}

Есть и другие, более медленные \nt{GABA}-нейроны, выходы которых садятся на отдельные входы получателя. Они работают как запрет~\eqref{eq:veto} из главы~\ref{sec:gaba}: закрывают один поток \nt{GLUT}-импульсов, не трогая остальных.

Деление на быстрые и медленные \nt{GABA}-нейроны --- старая картина, и она неполна. Сейчас в коре различают по меньшей мере три крупных класса~\cite{Tremblay2016}, и третий устроен неожиданно: он тормозит не \nt{GLUT}-нейроны, а другие \nt{GABA}-нейроны, прежде всего те, что закрывают входы~\cite{Pfeffer2013}. Торможение торможения --- двойное отрицание. Такой нейрон \emph{открывает} ворота, не посылая получателю ни одного возбуждающего импульса. Включают его сигналы из других областей коры и широковещательные каналы, так что он работает как разрешающий вход целого участка сети: пока он активен, запреты сняты, и поток проходит~\cite{Pi2013}. В приложении~\ref{sec:othermediators} этот приём назван растормаживанием.
\section{Экономный язык}

Последнее, что известно о языке нейронов твёрдо, --- он дорог. Импульс вместе с работой синапсов, которую он вызывает, обходится примерно в $10^9$ молекул АТФ (оценка для коры грызунов~\cite{Attwell2001}), а одна молекула даёт около $10^{-19}$ джоуля. Значит, импульс стоит порядка $E_{\text{имп}}\sim10^{-10}$ джоуля. Весь мозг потребляет около $20$ ватт, и если на передачу сигналов уходит половина, $P\sim10$ ватт, то средняя частота одного нейрона
\begin{empheq}[box=\keybox]{equation}
\bar r \;\approx\; \frac{P}{N\,E_{\text{имп}}}
\;\sim\; \frac{10\ \text{Вт}}{8.6\cdot10^{10}\cdot10^{-10}\ \text{Дж}}
\;\approx\; 1\ \text{импульс в секунду}.
\label{eq:energy}
\end{empheq}
Более подробные оценки для коры дают ещё меньше~\cite{Lennie2003}. Код мозга обязан быть \emph{разреженным}: быстро работать может лишь малая доля нейронов.

Из~\eqref{eq:spikeentropy} видно, что это не так уж плохо. При точности $1\ms$ импульс нейрона, работающего с частотой $100$ в секунду, несёт не больше пяти бит, а импульс нейрона, работающего раз в секунду, --- до одиннадцати. Если платить приходится за импульсы, выгодно говорить редко. Кора и правда меняет точность на энергию: при нехватке пищи нейроны сохраняют частоту импульсов, но тратят меньше на синаптические токи, и их ответы становятся менее точными~\cite{Padamsey2022}.

В азбуке Морзе частые буквы короткие: самая частая буква английского текста E --- одна точка. Код нейронов, насколько известно, следует тому же правилу в другой форме: то, что ожидаемо, стоит мало импульсов~\cite{Barlow1961}. Нейрон при постоянном раздражителе постепенно снижает частоту, датчики главы~\ref{sec:glutcomputer} отвечают на перемены, а не на уровни, а \nt{DOPA}-нейроны главы~\ref{sec:neurontypes} сообщают только об ошибке предсказания награды.

\begin{keyblock}
\begin{proposition}[Экономия]
\label{prop:economy}
Энергия ограничивает среднюю частоту нейрона величиной порядка одного импульса в секунду. Поэтому язык \nt{GLUT}-нейронов разрежен и тратит импульсы на новости: на то, что изменилось или не было предсказано. Энергетический предел измерен; что код мозга подстроен под наибольшее число бит на джоуль, --- хорошо обоснованная гипотеза.
\end{proposition}
\end{keyblock}

\section{Итог}

\begin{table}[t]
\centering
\footnotesize
\setlength{\tabcolsep}{4pt}
\renewcommand{\arraystretch}{1.15}
\begin{tabular}{>{\raggedright}p{2.6cm}>{\raggedright}p{2.3cm}>{\raggedright}p{3.0cm}>{\raggedright}p{2.0cm}>{\raggedright\arraybackslash}p{3.6cm}}
\toprule
Способ записи & Что несёт & Кто читает & Время на сообщение & Что известно \\
\midrule
номер сработавшего нейрона (код местом) & какое значение & любой получатель; детекторы совпадений & одна задержка & установлено у чувствительных нейронов и в определении направления на звук \\
частота одного нейрона & величину & интегратор с большой $\tau$; объём (гл.~\ref{sec:serocomputer}) & сотни мс -- секунды & установлено; для ступени обработки в $10$--$20\ms$ слишком медленно \\
частота группы & величину & нейрон с тысячами входов & $10$--$50\ms$ & установлено; выигрыш ограничен корреляциями~\eqref{eq:correlated} \\
время первого импульса, порядок & величину, порядок & детекторы совпадений с задержками & одна задержка & установлено на входе нервной системы; в коре обсуждается \\
фаза в местном ритме & величину, положение & нейроны с общим ритмом & цикл $12$--$50\ms$ и медленнее & наблюдается в отдельных областях; общая роль --- гипотеза \\
пачка («тире») & «важно»: надёжная доставка & облегчающийся синапс & $\sim10\ms$ & надёжность измерена; отдельный знак --- гипотеза \\
изменение частоты & новость & истощающийся синапс & $\sim0.1\,\text{с}$ & установлено \\
импульсы быстрых \nt{GABA}-нейронов & когда и насколько громко & все соседние \nt{GLUT}-нейроны & $1$--$30\ms$ & ритм и деление измерены; «пунктуация» как правило --- гипотеза \\
\bottomrule
\end{tabular}
\caption{Способы, которыми \nt{GLUT}- и \nt{GABA}-нейроны записывают сообщения импульсами. Правый столбец отделяет измеренное от предполагаемого.}
\label{tab:code}
\end{table}

Таблица~\ref{tab:code} собирает то, что мы знаем; видно из неё и то, чего мы не знаем. Мы не знаем, насколько кора пользуется точным временем импульсов, а не только частотами групп. Мы не знаем, есть ли у нейронов «слова» --- устойчивые сочетания импульсов многих нейронов, которые читаются как целое. И мы не знаем, одинаков ли язык в разных частях мозга. Азбуку Морзе можно выучить за вечер, потому что её таблица известна. Таблицу языка нейронов ещё предстоит составить, и составлять её будут, скорее всего, инженеры связи.

\section{Задачи}

\begin{exercise}[\lvlA]
\label{ex:04:1}
\emph{Ёмкость в числах.} $\Delta t=1\ms$. (а)~Сколько бит на джоуль даёт нейрон с частотой $1$ импульс в секунду при цене импульса из~\eqref{eq:energy}? (б)~При $r=10$ во сколько раз меньше границы~\eqref{eq:spikeentropy} извлекает получатель, который только считает импульсы за секунду?
\end{exercise}

\begin{exercise}[\lvlB]
\label{ex:04:2}
\emph{Оптимальная частота.} Нейрон тратит мощность $P_0+rE_{\text{имп}}$, где $P_0$ --- вторая половина из $20$~Вт, поделённая на все нейроны, $E_{\text{имп}}=10^{-10}$~Дж. При какой частоте $r$ число бит на джоуль $R_{\max}(r)/(P_0+rE_{\text{имп}})$ по~\eqref{eq:spikeentropy} с $\Delta t=1\ms$ наибольшее?
\end{exercise}

\begin{exercise}[\lvlB]
\label{ex:04:3}
\emph{Какая корреляция вредна.} $N=1000$ нейронов, дисперсия $\sigma^2=1$, попарная корреляция $c=0.1$. Вычислите информацию Фишера $\mathcal I=f'^{\mathsf T}\Sigma^{-1}f'$ ($\Sigma$ --- матрица ковариаций) для одинаковых наклонов $f'=\mathbf 1$ и для наклонов $\pm1$ с $\mathbf 1^{\mathsf T}f'=0$. Во сколько раз они различаются?
\end{exercise}

\begin{exercise}[\lvlB]
\label{ex:04:4}
\emph{Моделирование: детектор совпадений.} Нейрон~\eqref{eq:glutneuron} получает $1000$ пуассоновских входов по $5$ импульсов в секунду, $w=1$; порог $\theta$ таков, что свободное напряжение пересекает его раз в секунду. Моделью найдите, сколько одновременных входов $m$ зажигают нейрон с вероятностью $1/2$ при $\tau=10$ и $2\ms$.
\end{exercise}

\begin{exercise}[\lvlB]
\label{ex:04:5}
\emph{Проектирование: детектор относительных перемен.} Подберите синапс~\eqref{eq:depression}: установившийся ток при $40$ импульсах в секунду не больше чем на $10\,\%$ выше, чем при $10$, а после скачка до $40$ ток выходит на новый уровень с постоянной времени не больше $50\ms$. Каково наименьшее $\tau_{\text{вос}}$ и при каком $U$?
\end{exercise}

\begin{exercise}[\lvlA]
\label{ex:04:6}
\emph{Время первого импульса и пачки.} (а)~По~\eqref{eq:latency} найдите вход, при котором первый импульс приходит ровно через одну постоянную времени. От чего он зависит? (б)~При вероятности выброса $p=0.2$ сколько импульсов нужно в пачке, чтобы потерять все с вероятностью ниже $5\,\%$?
\end{exercise}

\begin{exercise}[\lvlC]
\label{ex:04:7}
\emph{Исследование: сколько бит в расположении.} Нейрон --- независимые клетки~\eqref{eq:spikeentropy}, $r=10$, $\Delta t=1\ms$. (а)~Разложите энтропию «слова» из $20$ клеток на энтропию числа импульсов и остаток. (б)~Смоделируйте оценку энтропии по частотам слов из $N=30$ и $10^4$ записей. Насколько она занижена?
\end{exercise}
